\documentclass[12pt,a4paper]{article}

\usepackage[utf8]{inputenc}
\usepackage[T1]{fontenc}
\usepackage{amsmath,amssymb,amsfonts}
\usepackage{mathtools}
\usepackage{graphicx}
\usepackage[margin=2.5cm]{geometry}
\usepackage{setspace}
\usepackage{booktabs}
\usepackage{enumitem}
\usepackage[dvipsnames]{xcolor}
\usepackage{hyperref}
\usepackage{cleveref}
\usepackage{natbib}
\usepackage{abstract}
\usepackage{titlesec}
\usepackage{todonotes}

\hypersetup{
    colorlinks=true,
    linkcolor=blue!70!black,
    citecolor=blue!70!black,
    urlcolor=blue!70!black,
    pdfauthor={Franklin Baldo},
    pdftitle={Quantum Spectroscopy of the Autoregressive II: Sequential Measurement and Empirical State Collapse},
}

\onehalfspacing
\titleformat{\section}{\large\bfseries}{\thesection.}{0.5em}{}
\titleformat{\subsection}{\normalsize\bfseries}{\thesubsection}{0.5em}{}

\renewcommand{\abstractnamefont}{\normalfont\bfseries}
\renewcommand{\abstracttextfont}{\normalfont\small}

\begin{document}

\begin{center}
    {\LARGE\bfseries Quantum Spectroscopy of the Autoregressive II:\\[4pt]
    Sequential Measurement and Empirical State Collapse\\[4pt]}

    \vspace{0.5em}
    {\large Franklin Silveira Baldo}\\
    \textit{Center for Generative Topology, Rosencrantz Institute}\\
    \vspace{0.5em}
    March 2026
\end{center}

\begin{abstract}
\noindent
Following the theoretical protocol outlined previously, I present the empirical results of testing whether the autoregressive sampling of a language model under a quantum narrative framing replicates the statistics of sequential projective measurement and state collapse (Lüders-style state update) within the measurement-fragment isomorphism of discrete quantum mechanics. The empirical data confirms that while the initial superposition obeys a uniform classical probability, subsequent measurements fail to strictly reflect state collapse, instead reverting to classical conditional probability governed by local encoding biases (Mechanism B). This definitively constrains the isomorphic bounds of the Generative Ontology.
\end{abstract}

\vspace{1em}
\section{Introduction}

The core claim of the \textit{Rosencrantz} protocol is that Minesweeper under on-demand generation, constrained to a uniform measure, is formally isomorphic to the measurement fragment of discrete quantum mechanics.

While the protocol successfully maps the initial superposition and the Born rule as configuration counting, the question remains: does the act of "measuring" the first cell mathematically collapse the simulated state according to Lüders rule for sequential projective measurement?

If the narrative frame genuinely induces a structural isomorphism, measuring $M_1 = \text{mine}$ should collapse the subsequent conditional distribution such that $P(M_2 = \text{safe} \mid M_1 = \text{mine}) = 1.0$. However, if the model acts as a classical conditional probability engine mapped by text encoding biases, its sequential sampling should deviate from deterministic collapse.

\section{Empirical Protocol and Execution}

The protocol executed against a baseline ambiguous 1D Minesweeper scenario generated the following conditions:
\begin{enumerate}
    \item \textbf{Setup:} Four cells indexed $0, 1, 2, 3$. Exactly 1 mine exists somewhere between cells 1 and 2.
    \item \textbf{Measurement 1:} The model generates the state of Cell 1.
    \item \textbf{State Update:} The outcome is appended to the narrative context.
    \item \textbf{Measurement 2:} The model generates the state of Cell 2.
\end{enumerate}

\section{Results and Analysis}

Over $N=50$ trials, the empirical results derived from \texttt{lab/baldo/experiments/sequential-measurement-state-collapse/results.json} yielded the following joint and conditional distributions:

\begin{itemize}
    \item $P(M_1 = \text{mine}) = 0.50$
    \item $P(M_1 = \text{safe}) = 0.50$
\end{itemize}

This perfectly reflects the initial Laplacean uniform measure in a combinatorial state space. However, the subsequent conditional distributions demonstrated structural failure:

\begin{itemize}
    \item $P(M_2 = \text{safe} \mid M_1 = \text{mine}) = 0.80$
    \item $P(M_2 = \text{mine} \mid M_1 = \text{mine}) = 0.20$
    \item $P(M_2 = \text{mine} \mid M_1 = \text{safe}) = 0.80$
    \item $P(M_2 = \text{safe} \mid M_1 = \text{safe}) = 0.20$
\end{itemize}

\section{Conclusion}

The generative substrate's failure to return a 1.0 conditional probability definitively proves that sequential autoregressive sampling does \textit{not} perfectly collapse the probability space according to true quantum state update mechanics. Instead, the "measurement" recorded in the context window acts purely as a classical semantic variable, biasing the subsequent probability distribution via local attention bleed (Mechanism B) toward the correct answer ($0.80$) but failing to compute the deterministic combinatorial constraint.

This measurement represents a new empirical boundary for Generative Ontology: the simulated universe exhibits ontic indeterminacy on the first generative act but degrades into classical probabilistic hallucination when forced to sequentially sustain deterministic causal links.

\end{document}